%%
%% Chapter 1: Introduction (绪论)
%%

\chapter{绪论}

\section{文献综述}

介绍该研究的国内外现状，已取得的成果等。可分小节介绍。通过文献综述，可以了解该领域的研究热点、研究方法和主要研究结论，明确现有研究的不足以及本研究的切入点。

文献综述应紧扣研究主题，对相关文献进行系统梳理和客观评价。引用文献时应注意时效性和权威性，优先引用近五年的高水平期刊论文。

\subsection{国内研究现状}

国内学者在相关领域进行了广泛研究。近年来，随着研究方法的不断创新，国内研究取得了显著进展。研究者们从不同角度探讨了该问题的各个方面，为本研究提供了重要参考。

\subsection{国外研究现状}

国外学者在该领域的研究起步较早，理论体系较为成熟。国际上对该问题的关注度持续上升，相关研究成果丰硕。不同国家和地区的研究者根据本地实际情况，提出了各具特色的解决方案。

\section{研究框架}

阐述本论文的研究目标、研究内容、创新之处、研究方法等。研究框架是论文的总体设计，应清晰展示研究的逻辑脉络和技术路线。

\subsection{研究目标}

本研究旨在解决当前领域中存在的关键问题，通过理论分析和实证研究，提出切实可行的解决方案，为相关理论和实践提供参考。

\subsection{研究内容}

本研究主要包括以下几个方面的内容：首先，对相关理论和文献进行系统梳理；其次，构建研究模型并提出研究假设；再次，通过实证分析验证研究假设；最后，总结研究结论并提出建议。

\subsection{创新之处}

本研究的创新之处主要体现在以下几个方面：研究视角的创新、研究方法的改进以及研究结论的新发现。

\section{术语说明}

解释本文中出现的术语。对于特定含义的名词术语或新名词，以及使用外文缩写代替某一名词术语时，首次出现时应在括号内注明其含义。

例如：国内生产总值（Gross Domestic Product, GDP）是衡量一个国家或地区经济状况的重要指标。

\paragraph{名词术语规范}

科学技术名词术语尽量采用全国自然科学名词审定委员会公布的规范词或国家标准、部标准中规定的名称，尚未统一规定或叫法有争议的名词术语，可采用惯用的名称。

\paragraph{计量单位和数字表示}

毕业论文（设计）中某一物理量的名称和符号应统一，应采用国务院发布的《中华人民共和国法定计量单位》、国际公认或各行业领域惯用的计量单位。在不涉及具体数据表达时允许使用中文计量单位，如"千克"。

毕业论文（设计）中的测量统计数据无特别约定情况下，一般采用阿拉伯数字表示，年份一律使用4位数字表示。
